\chapter{引用}
\label{ch:07}

按值函数可以计算新结果，但有些操作需要直接更新调用方的计数器，或只读取已有的大对象。引用把这种“访问同一个对象”的关系写进类型，同时要求来源在使用期间有效。

本章用一写一读两个函数说明 mut\& 与 const\&，再讨论来源、缺失和逃逸。引用的方便之处在于共享访问，其约束也来自同一个事实：它借用的数据并不归自己独立拥有。

\section{只读借用与可写借用}

上一章保留原值的方式很适合试算，但计数器加一常常希望修改原对象。我们需要在接口中表达这个区别：mut\& 借用可写对象，const\& 只借用读取能力。调用处的 \&value 指明从哪个对象取得引用。

下面先让 increment 修改 score，再让 read 读取它。两个函数围绕同一个整数工作，却有不同权限。只读引用不能用于修改对象的方法；可变引用也要求实际来源可写，不能靠形参上的 mut 凭空取得权限。

\par
\medskip
\begin{lstlisting}[style=chali]
void increment(int64 mut& value)
{
    value += 1;
}

int64 read(int64 const& value)
{
    return value;
}

void main()
{
    int64 score = 3;
    increment(&score);
    println("{}", read(&score));
}
\end{lstlisting}
\par
\medskip

\par
\medskip
\begin{lstlisting}[style=out]
4
\end{lstlisting}
\par
\medskip

\subsection*{代码详解}

increment 的形参带 mut\&，表示通过借用修改调用方对象；调用时的 \texttt{\&score} 明确提供 score 的来源。函数中 \texttt{value +=\allowbreak{} 1} 因而把 score 从 3 改为 4，而不是只修改一份整数副本。

read 的 const\& 只允许通过这个接口读取。它把读到的整数作为返回值交给 println，得到 4。两次借用按调用先后使用同一个仍然存活的 score；不要据此同时保留彼此冲突的可变访问。区别参数上的引用类型与调用处的取引用写法，才能判断函数究竟读写谁。

\section{引用需要一个来源}

引用在声明时绑定来源，没有以 null 表示的普通空引用。需要“可能借不到”时，使用 Option<T const\&> 或 Option<T mut\&>；None 表示没有引用，不构造空地址。

引用的副本仍借用同一对象。允许别名共享并不取消可写性和来源检查；本语言不采用全局独占可变借用规则。

\section{寿命、失效与逃逸}

返回引用时，来源需要在调用结束后继续存在。不能返回函数局部对象的引用，也不能让保存借用的闭包超过来源寿命。容器结构修改、资源释放和移动还可能使已有借用失效。

诊断 CHL1204 指向借用来源与使用路径。修正方法通常是返回拥有值、让调用方提供更长寿命的对象，或缩小借用范围。不要用裸地址转换掩盖来源问题。

\section{选择参数形式}

小型可复制值可按值传递；需要共享已有对象时选择引用，并写清是否修改。返回拥有值表示调用方接收资源，返回引用表示继续借用来源。第 8 章把这些来源规则与作用域清理连接起来。

\subsection*{小结}

把一条引用画成从使用位置指向来源的关系，再检查来源是否可写、是否仍存活、是否被移动或释放。普通引用必须有来源；需要表达缺失时让 Option 承担这个状态，而不是留一个未初始化引用。

本语言允许别名关系，但这不意味着任何重叠访问都合法。判断具体诊断要看访问方式和来源路径，不能照搬另一门语言的全局独占规则。接下来把来源检查与资源退出联系起来。

\section{习题}

\begin{enumerate}

\item 把 increment 参数改为按值，观察结果。

\item 能否将 const 变量传给 mut\&？

\item 如何表示可能不存在的引用？

\item 为什么局部引用不能返回？

\end{enumerate}

\subsection*{参考解答与验收要点}

\begin{enumerate}

\item 原 score 不因修改副本而增加。

\item 不能；修改接口或提供可写来源。

\item 使用对应 Option 引用类型并处理 None。

\item 函数退出时局部来源结束，返回值会借用已结束的对象。

\end{enumerate}

下一章：\hyperref[ch:08]{第 8 章}。

\section{番外：借来的荆州}

改写自《三国演义》鲁肃索荆州的相关情节。

却说鲁肃奉命讨还荆州，入府见关公，先递一纸：“某欲借阅库中账目。”关公捋髯道：“取只读通行牌来。”肃大喜，接牌便要改库房户名。

关公横刀拦住：“借你查阅，何曾许你过户？”肃曰：“既已取得引用，理当附送所有权。”关公道：“若依此言，你看我一眼，岂不连赤兔马也须骑走？”

鲁肃自觉理亏，又问：“倘若库房拆了，此牌可还能用？”关公笑道：“你拿旧门牌找新房子，倒比讨荆州还大胆。”

肃乃另取可写牌，关公将约定一一写明：所指何物，准做何事，何时到期。肃嫌文书太长，要在后面补“永久有效”，关公只将青龙偃月刀往桌上一放。肃立刻称赞这套生命周期检查十分直观。

席间关公持刀赴会，肃举杯未饮，又想把借牌交给下一任使者。关公道：“你得先看看来源还在不在。可写，也不是长生不老。”

鲁肃退回船上，随从问荆州何在。肃摊开手中通行牌，长叹：“今日只取得访问能力，连一块地砖也没有移进船里。”

\subsection*{技术对照}

引用表达访问关系而非自动拥有关系。只读引用不能用来写入；可写引用也不延长目标生命周期。返回或保存引用时要追踪真实来源。
